\documentclass[11pt,a4paper]{article}
\usepackage[utf8]{inputenc}
\usepackage[T1]{fontenc}
\usepackage{amsmath,amssymb}
\usepackage{graphicx}
\usepackage{geometry}
\usepackage{float}
\usepackage{xcolor}
\geometry{a4paper, margin=1in}

\title{Technical Report: Clinical vGFR Prediction Model (Round 3)}
\author{Antigravity Biomedical AI Assistant}
\date{March 25, 2026}

\begin{document}

\maketitle

\section{Establishing a classical machine learning model for clinical GFR data - v1}


This report details the machine learning model (\textit{Round 3 Champion}) designed to predict the corrected Glomerular Filtration Rate (\texttt{eGFRc}) from CT scan annotations. This model is optimized for small cohorts where traditional statistical methods might fail due to "overfitting" (finding patterns in noise).

\subsection{Cohort Description}
The analysis was performed on the \textbf{25-11-2025} cohort. 
\begin{itemize}
    \item \textbf{Sample Size ($n$):} 25 patients.
    \item \textbf{Data Availability:} Every patient has a complete set of three-phase CT scans (Arterial, Venous, and Equilibrium/Late phases) alongside demographic data (Age). All participants are male.
    \item \textbf{Clinical Labels:} The ground truth for this model is the corrected Estimated Glomerular Filtration Rate (\texttt{eGFRc}), which was calculated from serum creatinine measurements using standardized clinical formulas. It is considered the closest eGFR measurement from the scan date.
\end{itemize}


\subsection{Technical Foundations}
To understand how the model arrived at its predictions, we define four core modeling concepts used in this study.

\subsubsection{Ordinary Least Squares (OLS)}
\textbf{The Baseline:} OLS is the most common way to draw a "line of best fit" through data points. It tries to minimize the total squared distance between the actual measured values and the predicted values. 
\begin{itemize}
    \item \textit{The Limitation:} In small datasets like ours ($n=25$), OLS often performs "too well" on the training data by memorizing individual outliers, leading to poor results on any new patients.
\end{itemize}

\subsubsection{Regularization (The "Safety Brake")}
\textbf{The Control:} Regularization is a technique that adds a mathematical "penalty" to the model based on the size of its internal weights (coefficients). 
\begin{itemize}
    \item \textit{Why it's used:} It prevents the model from relying too heavily on any single feature or becoming overly sensitive to tiny variations. It forces the model to find the most stable, robust relationships that are likely to persist in future patients.
\end{itemize}

\subsubsection{Ridge Regression}
\textbf{The Algorithm:} Ridge Regression is a specialized version of OLS that incorporates \textbf{L2 Regularization}. 
\begin{itemize}
    \item \textit{The Mechanism:} Instead of just minimizing the distance to the points, Ridge also tries to keep the total "sum of squared weights" small. This creates a "shrunken" model that is much more reliable and less prone to erratic predictions when faced with a new patient.
\end{itemize}

\subsubsection{Stepwise Standard Selection}
\textbf{The Curator:} We have dozens of potential features, but using all of them ($\approx 65$ features for $25$ patients) is risky. Stepwise selection is a greedy search algorithm:
\begin{itemize}
     \item \textit{How it works:} It starts with zero features. It tests every feature and picks the one that single-handedly explains the most variance. Then it looks for a second feature that adds the most value \textit{on top of the first}, and so on. We set a limit of maximum 12 features, to further avoid over-fitting.
    
\end{itemize}

\subsection{Clinical Metrics and Stability}

\subsubsection{Leave-One-Out Cross-Validation (LOOCV)}
To simulate real-world usage, we used LOOCV. This means we trained the model 25 times; each time, we left exactly one patient out, trained on the other 24, and then "tested" the model by asking it to predict the missing patient. The metrics below represent these "blind" tests.

\subsubsection{Performance Results}
\begin{itemize}
    \item \textbf{Mean Absolute Error (MAE):} $9.71 \, \text{ml/min/1.73m}^2$. On average, our prediction is off by about 9 units.
    \item \textbf{$R^2$ (Signal Strength):} $0.289$. The model successfully explains ~29\% of the variation in kidney function purely from CT image data.
    \item \textbf{Stability ($p=0.001$):} We tested if this $R^2$ was due to "luck" by scrambling the data 1,000 times. The resulting p-value of 0.001 proves that the model is finding a genuine physiological signal.
\end{itemize}

\subsubsection{Model Weights and Feature Calculation}
The model intercept is \textbf{95.61}. This is the initial baseline prediction from which the model adds or subtracts based on the features. The "Weights" in the table below represent the \textbf{Direction and Magnitude of Influence}:

\begin{itemize}
    \item \textbf{Positive Weight ($+$):} A \textbf{directly proportional} relationship. As the measured value (e.g., volume) increases, the model's predicted vGFR increases (indicating better kidney function).
    \item \textbf{Negative Weight ($-$):} An \textbf{inversely proportional} relationship. As the measured value increases, the predicted vGFR decreases. This often happens in vascular features where high residual tracer in an artery might indicate it hasn't successfully moved into the filtration units yet at that specific phase of the scan.
\end{itemize}

\begin{table}[h]
\centering
\small
\begin{tabular}{llrl}
\hline
\textbf{Rank} & \textbf{Full Feature Name} & \textbf{Weight} & \textbf{Calculation / Measurement Method} \\
\hline
1 & \texttt{arterial\_venacava\_below\_kidney} & 1.25 & Point HU measurement in infra-renal IVC. \\
2 & \texttt{arterial\_right\_kidney\_vein} & 4.50 & Point HU in Inferior Vena Cava (IVC) below kidneys. \\
3 & \texttt{arterial\_venecava\_between\_kidney\_hepatic} & 2.72 & Point HU in Inferior Vena Cava (IVC) below kidneys. \\
4 & \texttt{arterial\_left\_kidney\_artery} & -1.31 & Point HU measurement in left renal artery. \\
5 & \texttt{venous\_right\_kidney\_artery} & 1.75 & Point HU (Venous phase) in right renal artery. \\
6 & \texttt{venous\_kidney\_hu\_mean} & -3.21 & \textbf{Mean} of all pixel HUs in the 3D kidney mask. \\
7 & \texttt{arterial\_left\_kidney\_vein} & 2.02 & Point HU measurement in left renal vein. \\
8 & \texttt{arterial\_aorta} & -1.72 & Point HU measurement in abdominal aorta. \\
9 & \texttt{late\_right\_kidney\_vein} & 3.14 & Point HU (Late phase) in right renal vein. \\
10 & \texttt{late\_kidney\_vol} & 2.18 & \textbf{Sum} of all voxels in 3D mask ($cm^3$). \\
11 & \texttt{arterial\_right\_kidney\_artery} & -2.37 & Point HU measurement in right renal artery. \\
12 & \texttt{venous\_venacava\_between\_kidney\_hepatic} & 1.85 & Point HU (Venous phase) in mid-renal IVC. \\
\hline
\end{tabular}
\caption{Model features, weights, and derivation logic.}
\end{table}

\begin{figure}[ht]
    \centering
    \includegraphics[width=0.85\textwidth]{figures/r3_feature_relevance.png}
    \caption{Feature Relevance: Standardized Ridge Coefficients.}
\end{figure}

\FloatBarrier






\subsection{Conclusion}
The Ridge-based approach provides a scientifically robust method for extracting vGFR from CT images. By using strict "safety brakes" (regularization), the model avoids fitting to noise and provides a reliable estimate for use in clinical research environments.

\begin{figure}[ht]
    \centering
    \includegraphics[width=0.7\textwidth]{figures/r3_champion_plot.png}
    \caption{Visual Correlation: Predicted vs. Actual Function.}
\end{figure}

\FloatBarrier
\end{document}
